% ECONOMICS_PROBLEM: {"id": "D83-243", "title": "Information design with endogenous acquisition", "jel_code": "D83", "source_id": "OP-0010"}
% !TeX program = xelatex
\documentclass[11pt,a4paper]{article}
\usepackage[margin=23mm]{geometry}
\usepackage{fontspec}
\setmainfont{Latin Modern Roman}
\usepackage{amsmath,amssymb,mathtools}
\usepackage{xcolor,enumitem,booktabs,longtable,array,xurl,hyperref}
\definecolor{muted}{HTML}{53636C}
\hypersetup{colorlinks=true,urlcolor=blue,linkcolor=blue}
\setlength{\parindent}{0pt}
\setlength{\parskip}{5pt}
\newcommand{\R}{\mathbb R}
\newcommand{\E}{\mathbb E}
\newcommand{\N}{\mathbb N}
\newcommand{\ind}{\mathbf 1}
\newcommand{\Var}{\operatorname{Var}}
\newcommand{\Cov}{\operatorname{Cov}}
\newcommand{\supp}{\operatorname{supp}}
\DeclareMathOperator*{\argmin}{arg\,min}
\DeclareMathOperator*{\argmax}{arg\,max}
\newcommand{\entrymeta}[3]{{\sffamily\small\color{muted}\textbf{\texttt{#1}}\quad|\quad #2\par #3\par}}
\newcommand{\Problemref}[1]{\texttt{#1}}
\newcommand{\JELref}[2]{\texttt{#2} (JEL \texttt{#1})}
\begin{document}
\section*{Information design with endogenous acquisition}
\textbf{Problem ID:} \texttt{D83-243}\quad \textbf{JEL:} \texttt{D83}\par
% BEGIN ECONOMICS STATEMENT
\label{problem:OP-0010}
\label{atlas:prob:M10}

\noindent\textbf{Original atlas ID:} M10. \textbf{Formulation:} editorial specification of a research family.

\subsection*{Definitions and assumptions}

Let the payoff state $\theta$ belong to finite $\Theta$, with full-support prior $\mu_0$. A sender commits publicly to an experiment $\kappa(m\mid\theta)$ on a specified finite message set $M$. After observing $m$, the receiver may select a further experiment $\xi(z\mid\theta,m)$ with finite signal alphabet $Z$ from a declared nonempty compact menu $\Xi_m$ of stochastic kernels, paying a finite nonnegative lower-semicontinuous cost $C_m(\xi)$, then choose $a\in A$, where $A$ is finite. Receiver and sender payoffs are bounded functions $u_R(a,\theta)$ and $u_S(a,\theta)$. Conditional posteriors follow Bayes' rule whenever the conditioning event has positive probability. Specify independence restrictions and what the receiver observes about the learning technology; zero-probability messages have no direct effect on expected payoffs. These menu and cost conditions ensure that a receiver optimum exists at every positive-probability message.

\subsection*{Formal research question}

Let $\mathcal B(\kappa,C,\Xi)$ be the receiver's optimal experiment-and-action policies, maximizing expected $u_R-C$ after every positive-probability message. Characterize
\[
\sup_\kappa\inf_{b\in\mathcal B(\kappa,C,\Xi)}
 \mathbb E_{\mu_0,\kappa,b}[u_S(a,\theta)].
\]
Determine which richer cost functions, multiple-receiver interactions, or sequential acquisition restrictions admit useful optimality and approximation results. For empirical design, characterize identification of $(C,\Xi)$ from observed acquisition choices and randomized prices or disclosures, and optimize over its identified set. Specify tie-breaking or pessimistic selection when the receiver has several optimal learning policies. A reduction to a distribution of posterior beliefs must prove that the relevant acquisition opportunities and costs are preserved.

\subsection*{Interpretation and scope}

Disclosure and subsequent private learning are jointly endogenous because the sender's message changes the receiver's value of information. An experiment that deters acquisition can be optimal in some cost classes, but that conclusion is not universal. Costly sender disclosure and costly receiver acquisition are separate mechanisms and should not be interchanged. Receiver learning that optimizes buyer surplus before monopoly pricing is also a distinct game. In richer settings, information supply can change other receivers' incentives through strategic interactions or public belief updates. Policy evaluation therefore needs the actual timing and commitment opportunities, and empirical cost estimates must distinguish technological access from willingness to pay for additional information.

\subsection*{Source audit and status}

Kamenica--Gentzkow (2011) and Roesler--Szentes (2017) are verified; original link [15] correctly identifies the former. Bergemann--Morris (2016) is ambiguous, so both relevant primary articles are listed. Crucially, Matysková--Montes (2023), added as [S4], already solves a sender-persuasion model with independent receiver learning and uniformly posterior-separable information costs. It establishes that broad endogenous-acquisition wording cannot itself be advertised as wholly unsolved. The objective above is editorial and contains solved subclasses. The prospective frontier must specify richer learning technologies, interactions, dynamic restrictions, or an empirical identification task that goes beyond those published results.

\subsection*{References}

\begin{enumerate}

\item[S1] Emir Kamenica and Matthew Gentzkow (2011). \emph{Bayesian Persuasion}. American Economic Review 101(6), 2590--2615. \url{https://www.aeaweb.org/articles?id=10.1257/aer.101.6.2590}. \textit{Locator:} abstract; sender-optimal signal characterization. \textit{Establishes:} Optimal persuasion with an exogenously specified receiver information environment. \textit{Does not establish:} general endogenous costly receiver learning.

\item[S2] Dirk Bergemann and Stephen Morris (2016). \emph{Bayes Correlated Equilibrium and the Comparison of Information Structures in Games}. Theoretical Economics 11(2), 487--522. \url{https://www.econtheory.org/ojs/index.php/te/article/viewArticle/20160487}. \textit{Locator:} abstract; obedience characterization. \textit{Establishes:} An outcome characterization and comparison of information structures using Bayes correlated equilibrium. \textit{Does not establish:} a general robust mechanism optimizer.

\item[S3] Anne-Katrin Roesler and Balázs Szentes (2017). \emph{Buyer-Optimal Learning and Monopoly Pricing}. American Economic Review 107(7), 2072--2080. \url{https://pubs.aeaweb.org/doi/10.1257/aer.20160145}. \textit{Locator:} abstract; bilateral trade model. \textit{Establishes:} Buyer-optimal information design in a monopoly-pricing model. \textit{Does not establish:} general sender persuasion with empirically unknown acquisition costs.

\item[S4] Ludmila Matysková and Alfonso Montes (2023). \emph{Bayesian Persuasion with Costly Information Acquisition}. Journal of Economic Theory 211, article 105678. \url{https://doi.org/10.1016/j.jet.2023.105678}. \textit{Locator:} abstract; uniformly posterior-separable cost model. \textit{Establishes:} A solution to a persuasion model with endogenous independent receiver learning. \textit{Does not establish:} a solution for unrestricted learning technologies or dynamic multi-receiver settings.

\item[S5] Dirk Bergemann and Stephen Morris (2016). \emph{Information Design, Bayesian Persuasion, and Bayes Correlated Equilibrium}. American Economic Review: Papers and Proceedings 106(5), 586--591. \url{https://pubs.aeaweb.org/doi/10.1257/aer.p20161046}. \textit{Locator:} abstract. \textit{Establishes:} Information-design interpretation of Bayes correlated equilibrium. \textit{Does not establish:} a resolution of endogenous costly learning in general.

\end{enumerate}

\subsection*{Common conventions: Source mathematical conventions}

Unless an entry states otherwise, agent and firm sets are finite. State and action spaces are measurable, strategies and outcome maps are measurable, probability laws are specified, and expectations exist for the targets under discussion. Infinite-horizon discount factors lie in $(0,1)$ and discounted objectives are finite. Norms, tolerances and complexity input sizes must be specified for computational comparisons. Constants or primitives that appear locally are inputs, not empirical facts asserted by this catalogue.

\textbf{Identification.} A statistical model $m\in\mathcal M$ implies an observable law $P_m$ and target $\theta(m)$. At law $P$, its identified set is
\[
\Theta(P;\mathcal M)=\{\theta(m):m\in\mathcal M,\ P_m=P\}.
\]
Point identification holds when this set is a singleton, or equivalently when $P_{m_1}=P_{m_2}$ implies $\theta(m_1)=\theta(m_2)$ for all compatible models. Sharp bounds mean bounds attained, or approached when the boundary is not attained, by compatible models. Writing this set is a definition, not a solution: the substantive task is to establish informative restrictions and derive or estimate the set. An unrestricted-confounding model may give only trivial bounds.

\textbf{Inference.} Identification, estimability and valid uncertainty quantification are separate questions. Fix $\alpha\in(0,1)$. A confidence procedure $C_n$ over a declared law class $\mathcal P$ has asymptotic uniform coverage at level $1-\alpha$ when
\[
\liminf_{n\to\infty}\inf_{P\in\mathcal P}\Pr_P\{\theta(P)\in C_n\}\geq1-\alpha.
\]
For set-identified targets the coverage event must be defined separately, for example coverage of the full identified set. If an entry uses identification-indexed models instead of uniquely indexed laws, the same coverage convention is applied over those models. Pointwise limits do not establish uniform validity, and asymptotic validity does not guarantee finite-sample precision. Sample sizes, dependence structures and weak-identification sequences are part of each application.

\textbf{Counterfactuals.} $Y_i(a)$ is a potential outcome under a specified intervention, including its timing, implementation and versions. It is not $Y_i$ conditional on voluntarily choosing $a$. A full intervention vector permits interference; shorthand scalar treatment notation does not silently impose no interference. Assignment, selection, attrition, anticipation and measurement must be specified. A claim about an instrument's local effect is not automatically a population policy effect.

\textbf{Economic equilibrium.} A counterfactual model includes best responses, constraints, feasibility, market clearing, state transitions and expectations consistent with the stated information. When several equilibria exist, either a declared selector is an input or the target is set-valued. Comparisons across policy regimes must use the stated selection convention. Reduced-form relationships are not presumed invariant after an equilibrium-changing intervention.

\textbf{Optimization and welfare.} Optimization is over the stated feasible set. If attainment is not established, $\sup$ or $\inf$ and $\varepsilon$-optimal choices replace an asserted maximizer or minimizer. For example, with value $V=\sup_{a\in\mathcal A}W(a)<\infty$, an $\varepsilon$-optimal set is $\{a\in\mathcal A:W(a)\geq V-\varepsilon\}$ for $\varepsilon>0$. Benefits, costs, risk and health or environmental outcomes can be aggregated only using declared units and conversion weights. Interpersonal utility weights, the treatment of fiscal transfers and foreign profits, discounting, and the behavioral welfare criterion are explicit inputs. A comparative-static derivative additionally requires existence and appropriate differentiability of the selected counterfactual.

\textbf{Sources.} Local labels [S1], [S2], and so forth restart in every question. Locators identify the relevant abstract, model, section or result at the level actually checked; a bibliographic or abstract-level verification is not represented as inspection of an entire proof. Working papers are distinguished from later publications and changing versions. Source summaries are paraphrased. All URL checks and editorial formulations refer to the audit date stated above.
% END ECONOMICS STATEMENT
\end{document}
